% =====================================================================
% 2.1 FUNDAMENTAÇÃO TEÓRICA
% =====================================================================
\subsection{Fundamentação teórica}
\label{subsec:fundamentacao}

\subsubsection{Timetabling educacional}

O termo \textit{timetabling} designa a atribuição de eventos a períodos e, conforme a formulação, a outros recursos, sujeita a restrições de viabilidade e a critérios de qualidade. \citeonline{schaerf1999survey} organiza os problemas educacionais em três categorias: \textit{school timetabling}, que atribui aulas periódicas a professores e turmas escolares; \textit{examination timetabling}, que agenda provas cujos conflitos decorrem dos estudantes inscritos; e \textit{university course timetabling}, que agenda aulas universitárias considerando professores, grupos de estudantes, períodos e salas.

As restrições exatas variam de uma instituição para outra. Por isso, o UCTP designa uma família de problemas, e não uma formulação única \cite{lewis2008survey,babaei2015survey}. Em geral, distinguem-se restrições fortes (\textit{hard}), que definem a viabilidade de uma grade, e restrições fracas (\textit{soft}), que medem sua qualidade e costumam compor uma função objetivo ponderada.

\subsubsection{Curriculum-Based Course Timetabling}

Duas formulações universitárias recorrentes diferem pela origem dos conflitos entre estudantes. No \textit{Post-Enrolment Course Timetabling}, os conflitos derivam das matrículas individuais. No \textit{Curriculum-Based Course Timetabling} (CB-CTT), eles são definidos por \textit{curricula}, isto é, grupos de disciplinas que um mesmo conjunto de alunos deve cursar em paralelo.

A trilha 3 da \textit{Second International Timetabling Competition} (ITC-2007) consolidou uma formulação de referência para o CB-CTT \cite{digaspero2007itc}. Nela, uma solução atribui a cada aula um período e uma sala. Indisponibilidades, conflitos de professor ou de currículo e a ocupação exclusiva de sala são restrições fortes. Capacidade da sala, número mínimo de dias de trabalho, compactação curricular e estabilidade de sala aparecem como critérios fracos. \citeonline{bettinelli2015overview} apresentam uma visão geral das formulações e dos métodos propostos para o CB-CTT.

Este trabalho usa o CB-CTT como referência, mas não o reproduz literalmente. No IC/UFF, a atribuição do professor é uma decisão do modelo, enquanto no benchmark da ITC-2007 o professor de cada disciplina é dado. Além disso, o horizonte é anual, a sala é decidida por encontro semanal e há regras locais de setores, horários fixos e carga docente.

\subsubsection{Complexidade computacional}

Uma versão simplificada do problema de horários contém o problema de coloração de vértices. Considere um grafo de conflitos $G=(V,E)$ no qual cada vértice representa um evento e cada aresta indica que dois eventos não podem ocorrer simultaneamente. Associar um de $K$ períodos a cada evento, sem conflitos, equivale a obter uma $K$-coloração de $G$, isto é,
\begin{equation}
  c(u)\neq c(v), \qquad \forall (u,v)\in E.
\end{equation}
Como decidir a existência de uma $K$-coloração é NP-completo para $K\geq 3$ \cite{garey1979computers}, as formulações de otimização que generalizam essa versão são NP-difíceis. Restrições adicionais de salas, recursos e preferências não eliminam essa dificuldade, o que motiva o uso de métodos heurísticos em instâncias de porte realista.

\subsubsection{Meta-heurísticas de trajetória}

Meta-heurísticas exploram o espaço de soluções sem garantir otimalidade, buscando soluções de boa qualidade dentro de um orçamento computacional \cite{talbi2009metaheuristics}. Seja $\mathcal{S}$ o espaço de soluções e $F:\mathcal{S}\rightarrow\mathbb{R}$ uma função a minimizar. Uma estrutura de vizinhança $N$ associa a cada solução $s$ o conjunto $N(s)$ de soluções alcançáveis por um movimento elementar.

O \textit{Simulated Annealing} (SA), proposto por \citeonline{kirkpatrick1983optimization}, aceita ocasionalmente movimentos de piora para reduzir o risco de estagnação em ótimos locais. Dada uma solução corrente $s$ e uma vizinha $s'$, seja $\Delta=F(s')-F(s)$. Movimentos com $\Delta\leq 0$ são sempre aceitos; quando $\Delta>0$, a aceitação ocorre com probabilidade
\begin{equation}
  P(\text{aceitação})=\exp\left(-\frac{\Delta}{T}\right),
  \label{eq:metropolis}
\end{equation}
em que $T>0$ é a temperatura. A temperatura é reduzida ao longo da execução, o que torna a aceitação de pioras cada vez menos provável.

O \textit{Iterated Local Search} (ILS) alterna busca local, perturbação da melhor solução e um critério de aceitação, explorando uma sequência de ótimos locais \cite{lourenco2003iterated}. O \textit{Variable Neighborhood Search} (VNS) explora sistematicamente mais de uma estrutura de vizinhança, alternando entre elas para combinar intensificação e diversificação \cite{hansen2001variable}.

\subsubsection{OptFrame e pyoptframe}

O OptFrame é um framework para o desenvolvimento de métodos de otimização combinatória baseado em componentes reutilizáveis \cite{coelho2010optframe,coelho2020microbenchmark}. Entre esses componentes estão a representação da solução, o método construtivo, o avaliador, os movimentos, as estruturas de vizinhança e as próprias meta-heurísticas. O pyoptframe fornece \textit{bindings} Python para o núcleo funcional do OptFrame \cite{pyoptframe2026}.

A separação entre componentes permite trocar a meta-heurística ou a vizinhança sem reescrever o modelo do problema. Essa propriedade motivou a escolha do framework: representação, avaliação e movimentos puderam ser implementados e testados em Python e, depois, registrados no motor nativo.

\subsubsection{Trabalhos relacionados}
\label{subsubsec:relacionados}

Os levantamentos de \citeonline{lewis2008survey} e \citeonline{babaei2015survey} mostram que meta-heurísticas de trajetória e populacionais são amplamente empregadas em problemas de timetabling universitário. A visão geral de \citeonline{bettinelli2015overview} reúne formulações e abordagens para o CB-CTT, incluindo métodos exatos e heurísticos avaliados sobre as instâncias da ITC-2007.

\pendente{completar com trabalhos aplicados a instituições brasileiras (por exemplo, publicações do SBPO) e com estudos que integram atribuição de professores e alocação de salas. Para cada trabalho, indicar o problema tratado, o método e a diferença em relação a este TCC: horizonte anual, duas grades com turmas compartilhadas, sala por encontro e carga mínima anual de obrigatórias.}
